\section{Agente de pronóstico}

Pronostica irradiancia y humedad de suelo para San Carlos. Es un agente LLM con
herramientas, pero el reparto de trabajo es estricto.

\begin{clave}[title=El LLM nunca toca los números]
Claude hace cuatro cosas: entiende la pregunta, traduce el horizonte a segundos, llama a
una herramienta que hace el cálculo físico, y redacta la respuesta con su incertidumbre.
Los números salen de un modelo con anclaje físico, no de la intuición del modelo de
lenguaje. Así la parte numérica es auditable y reproducible, y el LLM aporta lo que sí
sabe hacer: orquestación, ruteo y explicación.
\end{clave}

\begin{figure}[htbp]
\centering
\resizebox{\textwidth}{!}{% Figura: lazo "cerebro vs manos". El LLM orquesta; los numeros salen de las tools.
\begin{tikzpicture}[node distance=10mm]

\node[comp,text width=32mm] (q) {Pregunta en español\\\nota{``¿cuánta irradiancia}\\\nota{en 2 horas?''}};

\node[comp,text width=34mm,right=16mm of q] (llm) {\textbf{Claude}\\\nota{entiende, decide,}\\\nota{redacta}};

\node[comp,text width=38mm,right=18mm of llm] (tool) {\textbf{Herramienta}\\\nota{clear-sky + kt*}\\\nota{(persistencia inteligente)}};

\node[almacen,text width=28mm,right=14mm of tool] (db) {Store\\\nota{sólo lectura}};

\node[comp,text width=32mm,below=14mm of q] (r) {Respuesta + banda\\\nota{``$\approx$540 W/m$^2$}\\\nota{(410--660)''}};

\node[comp,text width=34mm,below=14mm of tool] (audit) {\cd{predicciones}\\\nota{write-back de auditoría}};

\draw[flecha] (q) -- (llm);
\draw[flecha] (llm) -- node[etiq,above]{\cd{forecast(\dots)}} (tool);
\draw[flecha] (tool) -- node[etiq,above]{SELECT} (db);
\draw[flecha] ([yshift=-4mm]tool.west) -- node[etiq,below]{\{valor, banda\}} ([yshift=-4mm]llm.east);
\draw[flecha] (llm) -- (r);
\draw[flechad] (tool) -- (audit);

\node[etiq,fill=none,below=2mm of llm,align=center]
  {el LLM \textbf{nunca} calcula:\\ todo número viene de una tool};

\end{tikzpicture}
}
\caption{El lazo de \emph{tool-use}. La flecha de vuelta lleva un diccionario con valor y
banda; el LLM sólo lo narra.}
\end{figure}

\subsection{La física}

El método es descomposición por cielo despejado. La idea: pronosticar la irradiancia
directa es difícil porque la domina la parábola diurna del sol, que es perfectamente
predecible. Si se divide por esa parábola, queda sólo el efecto de las nubes ---una señal
mucho más suave--- y \emph{eso} es lo que se persiste.

\begin{equation*}
  \mathrm{GHI_{medida}} = \mathrm{kt^*} \times \mathrm{GHI_{cielo\ despejado}}
\end{equation*}

La receta de \cd{smart\_persistence}, en tres pasos:

\begin{enumerate}[leftmargin=*,itemsep=2pt]
  \item Calcular kt* de los últimos 60 minutos, usando \emph{sólo} lecturas anteriores a
    \emph{ahora}.
  \item Tomar la \textbf{mediana} de esos kt* (robusta: un reflejo o una lectura atípica
    no arrastra el pronóstico).
  \item Reconstruir: $\mathrm{GHI_{pred}}(t+h) = \overline{\mathrm{kt^*}} \times
    \mathrm{GHI_{cs}}(t+h)$.
\end{enumerate}

El tercer paso mete la geometría solar del futuro, que es astronómica y por lo tanto
lícita: si $t+h$ cae más cerca del mediodía el pronóstico sube aunque las nubes no
cambien. Eso es exactamente lo que la persistencia ingenua ---repetir la última lectura---
no sabe hacer, y por eso se degrada en horizontes largos.

Para humedad de suelo no hay cielo despejado que descontar: se persiste la mediana
reciente, con banda por variabilidad. El suelo cambia lento y está muy autocorrelado.

\begin{aviso}[title=Barrera anti-fuga]
\cd{get\_recent\_data(now, lookback)} sólo devuelve lecturas con \cd{timestamp < now}.
Es una sola función, en una sola capa (\cd{data.py}), y es la que hace honesto al
\emph{backtest}: el reloj se simula y el pronosticador jamás ve el futuro que se le está
pidiendo predecir.
\end{aviso}

\subsection{Estructura y dirección de dependencias}

Las capas dependen sólo hacia abajo. Nadie importa hacia arriba:

\begin{center}
\cd{config} $\to$ \cd{domain} $\to$ \{\cd{data}, \cd{physics}\} $\to$
\cd{forecasters} $\to$ \cd{tools} $\to$ \cd{agent} $\to$ \cd{cli}
\end{center}

El valor de esa disciplina se ve al sustituir piezas: cambiar el sitio o el modelo es
tocar \cd{config}; agregar un pronosticador serio (AutoARIMA, ML) es implementar la
interfaz \cd{forecast(variable, horizon\_seconds, now)} sin tocar nada aguas arriba.

\subsection{La ingesta ambiental}

Un ETL propio, distinto del de CSV, trae datos de AgroDash al store:

\begin{itemize}[leftmargin=*,itemsep=2pt]
  \item \textbf{Fuente y destino son distintos a propósito}: \cd{DATABASE\_URL} apunta a
    AgroDash (sólo lectura) y \cd{STORE\_URL} a Supabase. El store desacopla el
    pronosticador de la fuente, le da historia propia y aloja predicciones y logs.
  \item \textbf{Idempotente e incremental}: \cd{ON CONFLICT (origen\_id) DO NOTHING},
    con marca de agua igual al máximo \cd{ts} del store menos un solape.
  \item \textbf{Escalable en memoria}: cursor del lado del servidor y COPY a tabla
    temporal en lotes de 50\,000. Los lotes chicos existen porque un COPY único de
    577\,000 filas chocaba contra el \cd{statement\_timeout} de Supabase.
  \item Los \emph{timestamps} \emph{naive} de AgroDash se etiquetan como hora local de
    Costa Rica, lo que se verificó viendo el pico diurno de irradiancia caer a las 11h.
\end{itemize}

\subsection{Endpoints}

\begin{center}
\small
\begin{tabular}{@{}p{40mm}p{78mm}p{14mm}@{}}
\toprule
\textbf{Ruta} & \textbf{Para qué} & \textbf{Clave} \\
\midrule
\cd{GET /health} & Ping para monitoreo. & no \\
\cd{GET /salud/ingesta} & Antigüedad de los datos por variable; responde \textbf{503} si
  están viejos, para que un \emph{healthcheck} lo note por código de estado. & no \\
\cd{GET /salud/panel} & Lo anterior más errores recientes, gasto del día y última
  predicción. & sí \\
\cd{POST /forecast} & Pronóstico a un horizonte; audita en \cd{predicciones}. & sí \\
\cd{POST /anomalias} & Detección determinista, sin LLM. & sí \\
\cd{POST /preguntar}, \cd{/chat} & Lazo completo del LLM, con traza. & sí \\
\cd{GET /serie}, \cd{/backtest}, \cd{/uso} & Datos y consumo. & sí \\
\bottomrule
\end{tabular}
\end{center}

Las dos rutas abiertas lo están por diseño: son de monitoreo automático y no exponen
datos de la serie, sólo edades y estado.

\subsection{Frenos de consumo}

Sin tope, un bucle o un \emph{scraper} dispara el costo del LLM. Hay tres frenos:
límite de ritmo por identidad en los endpoints que gastan tokens (12 por minuto),
límite más laxo en los deterministas (120 por minuto, que sólo protege CPU), y
presupuesto diario en dólares. El gasto se acumula en la tabla \cd{gasto\_diario} del
store, no en un archivo local: así sobrevive a que se recree el contenedor y no se
duplica si hay más de una instancia corriendo.

\subsection{Detección de anomalías}

Primer ladrillo de la futura capa de agentes, y deliberadamente sin LLM. Dado
(variable, ventana) devuelve hallazgos tipificados: \cd{sin\_datos\_recientes},
\cd{sensor\_plano} (el patrón de los 85\,°C), \cd{fuera\_de\_rango}, \cd{outlier}
(puntuación robusta por mediana y MAD) y \cd{drift}. La señal analizada es kt* para
irradiancia ---reutiliza el mismo clear-sky--- y el crudo para humedad. El LLM lo llama
como herramienta y sólo narra los hallazgos.
